\section{Lean 형식화}

2026-08-21.

Lean은 정리 증명기이자 함수형 프로그래밍 언어다. 정의·정리·증명을 기계가
한 단계씩 검사할 수 있는 형태로 다시 쓰는 것이 형식화다.

\begin{itemize}
\tightlist
\item
  기반: 귀납적 구성의 계산(Calculus of Inductive Constructions). 명제는
  타입, 증명은 그 타입의 항(커리--하워드 대응).
\item
  사례: Mathlib(공동 수학 라이브러리), Liquid Tensor Experiment, PFR
  추측의 형식화.
\item
  CS 쪽 연결: 의존 타입, 증명 보조기, 형식 검증(CompCert 컴파일러, seL4
  커널).
\end{itemize}

\subsection{직접 한 것}

\begin{itemize}
\tightlist
\item
  원 밖의 점에서 그은 접선 공식 \(r(px+qy)\pm\ell(py-qx)=rd^2\)를 Lean 4
  (v4.23.0, core only, Mathlib 없이)로 임의의 가환환 위에서 검증.
  회전·복소수 해석도 별도 파일로 검증. → \url{analytic-geometry.md}
\end{itemize}

\subsection{파일}

\begin{itemize}
\tightlist
\item
  \href{../lean/Tangent.lean}{\texttt{lean/Tangent.lean}}: 법선벡터
  \(n=(rp-s\ell q,\ rq+s\ell p)\)로 쓴 공식과 같음, 직선이 \(P\)를 지남,
  \(|n|^2=d^4\)(거리 \(=r\)), 원과의 공통점은 \(T=rn/d^2\) 하나뿐, 두
  접선(\(s=\pm1\))이 다름
\item
  \href{../lean/TangentSimple.lean}{\texttt{lean/TangentSimple.lean}}:
  \(P\)를 각 \(\alpha\)(\(\cos\alpha=r/d\), \(\sin\alpha=\ell/d\))만큼
  돌린 회전 형태, 복소수 형태 \(T\bar P=r(r+is\ell)\)
\end{itemize}

모두 Mathlib 없이 \texttt{grind} 전술로 증명하고, 정수 예제는
\texttt{decide}로 확인한다. 실행: Lean 4 v4.23.0에서
\texttt{lean\ lean/Tangent.lean}.
